% !PRINTED 68
\weberpage{68}
Der Inbegriff aller rationalen Functionen von $r$ mit rationalen Zahlen als Co\"efficienten ist also ein Zahlkörper $\Omega(r)$ vom Grade $\nu$, den wir einen \gesperrt{Kreistheilungskörper} nennen. Wir wollen aber den Begriff des Kreistheilungskörpers noch etwas allgemeiner fassen und darunter jeden Körper verstehen, dessen Zahlen lauter rationale Functionen irgend welcher Einheitswurzeln sind.

Den Körper $\nu^{\text{ten}}$ Grades $\Omega(r)$, der aus \gesperrt{allen} rationalen Functionen einer $m^{\text{ten}}$ Einheitswurzel $r$ besteht, nennen wir zur genaueren Unterscheidung den \gesperrt{vollen Kreistheilungskörper} der Ordnung $m$ und bezeichnen ihn mit $\Omega_m$.

Beliebige Einheitswurzeln $r, r', r'' \ldots$ beliebiger Grade $m, m', m'' \ldots$ kann man immer auffassen als Potenzen einer und derselben Einheitswurzel $\varrho$, deren Grad das kleinste gemeinschaftliche Vielfache von $m, m', m'' \ldots$ ist. Demnach ist ein Kreistheilungskörper, der nur rationale Functionen von $r, r', r'' \ldots$ enthält, ein Theiler des vollen Kreistheilungskörpers $\Omega(\varrho)$, und wir bekommen also alle überhaupt existirenden Kreistheilungskörper, wenn wir die sämmtlichen Divisoren aller vollen Kreistheilungskörper aufsuchen.

Die \gesperrt{Galois'sche} Gruppe des Körpers $\Omega_m$ besteht aus den sämmtlichen Substitutionen
\weq{3}{(r, r^n),}
wenn $n$ jede nach dem Modul $m$ genommene relative Primzahl zu $m$ bedeutet. Denn die Kreistheilungsgleichung $\nu^{\text{ten}}$ Grades, deren Wurzeln die $r^n$ sind, ist eine Normalgleichung und also ihre eigene \gesperrt{Galois'sche} Resolvente. Da, wenn $a, b$ zwei dieser Zahlen $n$ sind,
\[
(r, r^a)(r, r^b) = (r, r^{ab})
\]
ist, so ist diese Gruppe isomorph mit der Gruppe $\mathfrak{N}$ aller Zahlclassen $N$ der zu $m$ theilerfremden Zahlen, die wir im vorigen Paragraphen betrachtet haben.

Ist $\mathfrak{A}$ ein Theiler von $\mathfrak{N}$, und $\varrho$ eine zu $\mathfrak{A}$ gehörige Function aus $\Omega_m$, so ist der Inbegriff der rationalen Functionen von $\varrho$, $\Omega(\varrho)$, ein in $\Omega_m$ enthaltener Körper.

Ist umgekehrt $\Omega$ ein Theiler von $\Omega_m$ und $\varrho$ eine primitive Zahl des Körpers $\Omega$, also auch eine Zahl in $\Omega_m$, so kann $\Omega$ als der Inbegriff der rationalen Functionen von $\varrho$ dargestellt und mit $\Omega(\varrho)$ bezeichnet werden.
